\section{Results}
\label{sec:results}

\subsection{Candidate Disposition and Human Judgments}

The 44 packages contributed 440 candidates: 115 from 2018, 79 from 2019, 113 from 2024, and 133 from 2025 (Table~\ref{tab:cohort}). Human review retained 413 (93.9\%) and excluded 27 (6.1\%), with no deferred or incomplete entries. These totals begin at candidate registration; observations screened out earlier are outside this denominator. Final human domains comprised 173 statistical, 12 report-internal, 216 protocol--report, and 39 planning-material candidates. Exclusions comprised 18 statistical, 2 protocol--report, and 7 planning-material candidates. Five domain assignments changed during review.

Of 155 retained statistical candidates, 139 were confirmed, 15 probable, and 1 indeterminate. The confirmed group contained 37 substantive and 102 presentation-only issues. The probable group contained 4 substantive and 9 presentation-only issues, with substantiveness indeterminate for 2. The remaining indeterminate issue was classified as presentation only. Optional downstream-impact judgments were direct in 29 cases, possible in 40, unlikely in 85, and unknown in 1; these were assessments of potential use rather than observed consequences.

All 12 report-internal candidates were retained: 5 substantive, 6 minor, and 1 indeterminate. One substantive candidate had an original protocol--report relationship and was excluded from the temporal endpoint. Of 32 retained planning-material findings, 4 were substantive, 25 minor, and 3 indeterminate. The 214 retained protocol--report findings comprised 179 substantive differences, 14 minor differences, 6 version or timing ambiguities, and 10 indeterminate comparisons; 5 had no optional interpretation.

For illustration, a candidate concerning a binary risk-ratio identity was retained as a confirmed substantive statistical issue. A difference between planned and reported medication composites was retained descriptively even though the report disclosed that statin data had not been collected. A candidate concerning secondary-outcome and safety-analysis labels was excluded because those descriptions could coexist in the supplied plans. The supplement provides identifiers and further adjudication details.

\begin{table}[tbp]
\centering\small
\caption{Availability sample and final human dispositions.}
\label{tab:cohort}
\begin{tabular}{lrrrrr}
\toprule
Year & Papers & PDF pages & Candidates & Retained & Excluded \\
\midrule
2018 & 10 & 1070 & 115 & 107 & 8 \\
2019 & 9 & 994 & 79 & 76 & 3 \\
2024 & 10 & 1097 & 113 & 107 & 6 \\
2025 & 15 & 2398 & 133 & 123 & 10 \\
\midrule
Total & 44 & 5559 & 440 & 413 & 27 \\
\bottomrule
\end{tabular}
\par\smallskip\begin{minipage}{0.97\linewidth}\footnotesize The packages contained 188 PDFs and one supplementary Excel worksheet. Retention includes descriptive and uncertain findings; it does not imply confirmation.\end{minipage}
\end{table}


\subsection{Paper-Level Outcomes and Period Comparison}

The principal endpoint comprised 41 substantive statistical issues (37 confirmed and 4 probable) in 22 of 44 papers: 50.0\% (exact binomial 95\% confidence interval, 34.6\% to 65.4\%). By year, the eligible counts were 14, 6, 11, and 10. These are findings in the selected packages, not an estimate of journal-wide prevalence.

The earlier period had 20 eligible issues in 19 papers (mean, 1.05), compared with 21 in 25 later papers (mean, 0.84). The later-minus-earlier difference was $-0.21$ issues per paper (95\% bootstrap confidence interval, $-1.02$ to 0.56; exact permutation $P=0.654$). Papers with at least one eligible issue numbered 10/19 (52.6\%) and 12/25 (48.0\%), respectively ($P=1.000$). The intervals were compatible with differences in either direction.

Restricting the endpoint to confirmed substantive statistical issues gave 16 earlier and 21 later issues, or 0.84 per paper in both periods ($P=1.000$). All four probable substantive issues were from the earlier period. The secondary endpoint contained one earlier and three later report-internal inconsistencies ($P=0.622$); combining the two endpoints gave mean counts of 1.11 and 0.96 ($P=0.754$). Table~\ref{tab:temporal} presents the comparisons. Supplementary Table S1 gives all paper-level counts, and Supplementary Figure S1 displays the statistical counts by period.

\begin{table}[tbp]
\centering\small
\caption{Exploratory comparisons using exact paper-level permutation tests.}
\label{tab:temporal}
\begin{tabular}{>{\raggedright\arraybackslash}p{4.7cm}rrrr}
\toprule
Outcome & 2018--2019 & 2024--2025 & Difference (95\% CI) & $P$ \\
\midrule
Substantive statistical: count & 1.05 & 0.84 & -0.21 (-1.02, 0.56) & 0.654 \\
Substantive statistical: any & 10/19 & 12/25 & -4.6 (-33.7, 25.7) & 1.000 \\
Substantive report inconsistency & 0.05 & 0.12 & 0.07 (-0.09, 0.24) & 0.622 \\
Confirmed statistical only & 0.84 & 0.84 & 0.00 (-0.77, 0.74) & 1.000 \\
Combined substantive report & 1.11 & 0.96 & -0.15 (-1.02, 0.68) & 0.754 \\
\bottomrule
\end{tabular}
\par\smallskip\begin{minipage}{0.97\linewidth}\footnotesize Count rows show mean issues per paper; the binary row shows papers with at least one issue. Differences are later minus earlier, in counts or percentage points. Intervals use 20,000 within-period paper bootstrap resamples. Permutation $P$ values are two-sided, exact, and unadjusted. The last two rows are sensitivity analyses.\end{minipage}
\end{table}

